\chapter{Antecedentes}
%En esta sección se va desde aspectos generales a  aspectos específicos (como un embudo). No se olvide que es la primera parte que tiene contacto con el lector y que hará que este se interese en el tema a investigar.
%El objetivo de esta sección es llevar al lector hacie el tema que se va a tratar en forma específica y dejar la puerta abierta a otras investigaciones.


\section{Trabajos Relacionados}

Antecedentes o esbozo del estado del arte. Reflejan los avances y el Estado actual del conocimiento en un área determinada y sirve de modelo o ejemplo para futuras investigaciones. Se refieren a todos los trabajos de investigación que anteceden a nuestro, es decir, aquellos trabajos donde se hayan manejado las mismas variables o se hallan comparaciones y tener ideas sobre cómo se trató el problema en esa oportunidad (dos a cuatro páginas - según guía) (más información ver guía).


Existen dos tipos de citaciones: citep y citet, los ejemplos a continuación respectivamente: \cite{lee2007trajectory} y \cite{roh2010nncluster}.
